\documentclass[a4paper,10pt,twoside,openany]{article}

\usepackage[lang=hebrew]{maths}
\usepackage{hebrewdoc}
\usepackage{stylish}
\usepackage{lipsum}
\let\bs\blacksquare

\usepackage{siunitx,array}
\sisetup{table-format=-1.5, group-digits=false}
\newcolumntype{L}{>{\phantom{$-$}}l}       % prefix some whitespace
\newcommand\mcL[1]{\multicolumn{1}{L}{#1}} % handy shortcut macro

\setlength{\parindent}{0pt}

%%%%%%%%%%%%
% Styling %
%%%%%%%%%%%%

\usepackage{enumitem}

%%%%%%%%%%%%%
% Counters  %
%%%%%%%%%%%%%

\setcounter{section}{1}     
            
%BIBLIOGRAPHY
\usepackage[
backend=biber,
style=alphabetic,
]{biblatex}
\addbibresource{bibliography.bib} %Imports bibliography file

%%%%%%%%%%
% Title  %
%%%%%%%%%%
\title{
אלגברה ב' - גיליון תרגילי בית 7 \\
צורת ז'ורדן
\\
\vspace{1cm}
\large{תאריך הגשה: 18.6.2026}
}
\date{}

\begin{document}
\maketitle

\begin{exercise} %0
יהי
$\FF$
שדה ויהיו
$\lambda \in \FF, n \in \NN_+$.
\begin{enumerate}
\item
 הראו כי
\[\text{.} m_{J_n\prs{\lambda}} = \prs{x - \lambda}^n\]

\item הסיקו כי אם
$T \in \End_\FF\prs{V}$
אופרטור שהפולינום האופייני שלו מתפצל מעל
$\FF$,
ושהערכים העצמיים השונים שלו הם
$\lambda_1, \ldots, \lambda_m$,
מתקיים
\begin{align*}
m_T\prs{x} = \prod_{i \in [m]} \prs{x - \lambda_i}^{k_i}
\end{align*}
כאשר
$k_i$
גודל בלוק ז'ורדן המקסימלי עם ערך עצמי
$\lambda_i$
בצורת ז'ורדן של
$T$.
\end{enumerate}
\end{exercise}

\begin{exercise} %1
יהי
$n \in \NN_+$
ותהיינה
$A,B \in \Mat_n\prs{\CC}$
עבורן
\begin{align*}
p_A\prs{x} &= p_B\prs{x} \\
\text{.} m_A\prs{x} &= m_B\prs{x}
\end{align*}
הוכיחו או הפריכו את הטענות הבאות. עבור אלו שאינן נכונות, מיצאו דוגמה נגדית.

\begin{enumerate}
\item
$A,B$
דומות אם
$n=3$.

\item
$A,B$
דומות אם
$n=4$.

\item
$A,B$
דומות אם
$n=5$.
\end{enumerate}
\end{exercise}

\begin{exercise}%2
\begin{enumerate}
\item מצאו את צורת ז'ורדן של
$J_n\prs{\lambda}^{-1}$
עבור
$\lambda \in \mbb{C}\setminus\set{0}$
ועבור
$n \in \mbb{N}_+$.

\item תהי
$A \in M_n\prs{\mbb{C}}$
הפיכה.
מצאו את צורת ז'ורדן של
$A^{-1}$
כתלות בזאת של
$A$.

\item מצאו תנאי הכרחי ומספיק על הבלוקים בצורת ז'ורדן של
$A \in M_n\prs{\mbb{C}}$
כך שיתקיים
$A \sim A^{-1}$.
\end{enumerate}
\end{exercise}

\begin{exercise}
תהיינה
\begin{align*}
A_1 &\coloneqq \pmat{1 & -2 & 2 \\ 2 & -1 & 2 \\ 2 & -2 & 3}, \quad 
A_2 \coloneqq \pmat{1 & 2 & 2 \\ 2 & 1 & 2 \\ 2 & 2 & 3}, \quad
A_3 \coloneqq \pmat{1 & 2 & -2 \\ 2 & 1 & -2 \\ 2 & 2 & -3}
\end{align*}
מטריצות ב־%
$\Mat_3\prs{\RR}$,
ויהי
$v_0 = \pmat{3 \\ 4 \\ 5}$.

נגיד כי וקטור
$v = \pmat{a \\ b \\ c} \in \NN_+^3$
הוא
\emph{שלשה פיתגורית}
אם
\[\text{,} f\prs{v} \coloneqq a^2 + b^2 - c^2 = 0\]
וכי וקטור כזה הוא
\emph{שלשה פיתגורית פרימיטיבית}
אם בנוסף המחלק המשותף המקסימלי של
$a,b,c$
הוא
$1$,
ו־%
$b$
זוגי. (יכולנו לחלופין לדרוש כי $a$ זוגי. בדיוק אחד מהם זוגי כי אחרת הסכום $a^2 + b^2$ הינו זוגי אך לא מתחלק ב־%
$4$
ואז אין
$c \in \NN_+$
עבורו
$c^2 = a^2 + b^2$).

\begin{enumerate}
\item מתקיים כי אם
$v \in \RR^3$
שלשה פיתגורית, גם
$A_i v$
שלשה פיתגורית, וגם כי אם
$v \in \RR^3$
פרימיטיבית, גם
$A_i v$
פרימיטיבית.
הראו זאת עבור
$i = 1$
והסיקו כי
$A_1^k v_0$
שלשה פיתגורית פרימיטיבית לכל
$k \in \NN$.

\textbf{רמז:}
עבור פרימיטיביות, הראו שמקדמי המטריצה
$A_1^{-1}$
שלמים (אין צורך לחשב אותה לשם כך), והיעזרו בה.

\item ידוע כי $1$ ערך עצמי של
$A_1$
ושל
$A_3$,
וכי
$-1$
ערך עצמי של
$A_2$.

הראו כי
$A_2, A_3$
לכסינות, ומיצאו מטריצה
$P \in \Mat_3\prs{\RR}$
הפיכה, ומטריצת ז'ורדן
$J \in \Mat_3\prs{\RR}$
כך שמתקיים
$P A_1 P = J$.

\item ידוע כי כל שלשה פיתגורית פרימיטיבת מתקבלת על ידי כפל משמאל של
$v_0$
במטריצות מבין
$A_1, A_2, A_3$.
כלומר, קבוצת השלשות הפיתרגוריות הפרימיטיביות היא
\[\text{.} \set{\prs{\prod_{j \in [k]} A_{i_j}} v_0}{\substack{k \in \NN \\ \forall j \in [k]: i_j \in [3]}}\]

הראו כי הקבוצה
$\set{A_1^k v_0}{k \in \NN}$
היא בדיוק קבוצת השלשות הפיתגוריות הפרימיטיביות
$\pmat{a \\ b \\ c}$
עבורן
$c = b+1$,
וחשבו את
$A_1^k v_0$
כדי למצוא קבוצה זאת.
\end{enumerate}
\end{exercise}

\begin{remark}
מצאנו בעצם בתרגיל את כל המשולשים הישרי־זווית עם צלעות מאורכים שלמים, כך שאורך היתר גדול ב־%
$1$
מאורך אחת הצלעות.
\end{remark}

\begin{remark}
כדי להראות שקבוצת השלשות הפיתרוגיות הפרימיטיביות היא זאת שכתבנו, ניתן להראות כי עבור שלשה פרימיטיבית
$v = \pmat{a \\ b \\ c} \neq v_0$,
בדיוק אחד מבין הוקטורים
$A_i^{-1} v$
הוא בעל מקדמים שלמים, וכי הוא מהווה שלשה פיתגורית פרימיטיבית.

ניתן להראות זאת כתרגיל. את האי־שוויון $a+b < \frac{3}{2}c$ שמופיע כחלק מכך אפשר להראות ישירות, אך גם נדע להראות אותו בקלות בהמשך בעזרת מכפלות פנימיות.
\end{remark}

\end{document}